{\color{gray}\hrule}
\begin{center}
\section{Data and Code Availability}  % draft: restore \section{Conclusions} and the subsection heading below once the conclusions exist
\label{sec:availability}
\bigskip
\end{center}
{\color{gray}\hrule}
\vspace{0.5cm}
% TODO The answer to the research question in two or three sentences; the contributions restated as findings; future work that follows from the limitations. Back matter: data and code availability, references.
\iffalse
\placeholder{15}
\fi

% \subsection{Data and Code Availability}
The code is at \url{https://github.com/sclfnc/hypergraph-walks-bach-palestrina}. The data are not redistributed: the notebooks download every source. Both corpora are distributed with music21 \cite{cuthbert2010, cuthbert2026music21}, whose software licence, BSD 3-Clause, does not cover the files of its corpus. The Bach files state no licence, except three \textcopyright{} CCARH (1994). The Palestrina files are \textcopyright{} John Miller (2000). Both of these copyrights reserve the rights to derivative electronic formats. When in Rome \cite{gotham2023, gotham2026wir} is under CC BY-SA 4.0 and the Bach Chorales Melody-Harmony corpus \cite{sears2023} under CC BY 4.0.
